% Cuerpo editor-ready del nuevo capitulo 57.

\section{Sum--product defect of APP and linear coefficient of the modular product Euler characteristic}
\label{sec:l9s102:euler-defecto}

\subsection{Exact identity \texorpdfstring{\(459-405=54=[q]t_3(q)\)}{459-405=54=[q]t3(q)}}

On the support of \(81\) cells, the two APP leaves produce
\begin{equation}
 S_+=405,
 \qquad
 S_\times=459,
 \qquad
 S_\times-S_+=54.
 \label{eq:l9s102:euler-405-459}
\end{equation}
The triadic regularization supplies the exponent \(-1/12\). Its
dodecaphase publication determines the pole \(q^{-1}\), and the Dedekind quotient
of level \(3\) is written
\begin{equation}
 t_3(q)
 =\left(\frac{\eta(\tau)}{\eta(3\tau)}\right)^{12}
 =q^{-1}\prod_{n\ge1}(1+q^n+q^{2n})^{-12},
 \qquad q=e^{2\pi i\tau}.
 \label{eq:l9s102:euler-t3}
\end{equation}

\begin{theorem}[Euler--HMT sum--product identity]
\label{thm:l9s102:euler-54}
The defect between the two APP evaluations coincides with the linear
coefficient of \(t_3\):
\begin{equation}
 \boxed{
 S_\times-S_+=54=[q]\,t_3(q).}
 \label{eq:l9s102:euler-54}
\end{equation}
\end{theorem}

\begin{proof}
Each factor admits
\((1+q^n+q^{2n})^{-12}=1-12q^n+O(q^{2n})\). For the coefficient of
\(q\) subsequent to the pole \(q^{-1}\), the degree-\(2\) partitions
of the level-\(3\) product contribute. The certified finite expansion produces
\([q]t_3=54\). APP obtains the same integer by the exact subtraction
\(459-405\).
\end{proof}

\subsection{Conservation of the defect via residue, quotient, and nonary carry}

The identity \cref{eq:l9s102:euler-54} links the additive difference of two finite sheets
to an infinite modular product. The link preserves the origin of
\(54\): sum, product, and regularization remain three distinguishable
state coordinates.

\subsubsection{EMRH Completion and Conservation of the Unit}
\label{sec:l9s102:euler-emrh}

The ternary fiber of depth \(6\) has \(3^6=729\) elements. Its
decimal completion distinguishes
\begin{equation}
 729+270=999,
 \qquad
 729+271=1000,
 \qquad
 271=243+27+1.
 \label{eq:l9s102:euler-729-1000}
\end{equation}
The terminal term \(+1\) is the closure unit transmitted by the
carry. The binomial identity
\begin{equation}
 1000=(9+1)^3=9^3+3\cdot9^2+3\cdot9+1
 \label{eq:l9s102:euler-binomial-1000}
\end{equation}
expresses transport between the nonadic chart and decimal publication
without losing the genealogy of the completion.

Archimedean evaluation is performed on compatible half-closed cylinders.
The two expansions of a terminating rational represent the
same value and preserve distinct boundary histories. This structure
types the identity \(0.999\ldots=1\) as equality of evaluation accompanied
by route memory.

\section{Quadratic algebras of TRIT and the elliptic, parabolic, and hyperbolic exponential law}
\label{sec:l9s102:euler-algebras-trit}

\subsection{Operator family \texorpdfstring{\(J_{+1}^2=-I\), \(J_0^2=0\), and \(J_{-1}^2=I\)}{J+1^2=-I, J0^2=0, and J-1^2=I}}

Para \(\tau\in\{+1,0,-1\}\), definimos
\begin{equation}
 \mathbb A_\tau
 =\mathbb R[J_\tau]/(J_\tau^2+\tau I).
 \label{eq:l9s102:euler-algebras}
\end{equation}
The exponential series converges in each algebra and reduces to
\begin{align}
 e^{tJ_{+1}}&=\cos t\,I+\sin t\,J_{+1},
 \label{eq:l9s102:euler-eliptica}\\
 e^{tJ_0}&=I+tJ_0,
 \label{eq:l9s102:euler-parabolica}\\
 e^{tJ_{-1}}&=\cosh t\,I+\sinh t\,J_{-1}.
 \label{eq:l9s102:euler-hiperbolica}
\end{align}
TRIT assigns elliptic closure, first-order variation, and
hyperbolic opening, respectively. The classical identity occupies the leaf \(\tau=+1\):
\begin{equation}
 \boxed{\operatorname{Exp}(\pi J_{+1})+I=0.}
 \label{eq:l9s102:euler-clasica}
\end{equation}
The parabolic sheet preserves the first jet, and the hyperbolic sheet publishes
a pair of reciprocal eigenvalues.

\subsection{Identities \texorpdfstring{\(\operatorname{Exp}(\pi J_{+1})+I=0\)}{Exp(pi J+1)+I=0} and \texorpdfstring{\(\operatorname{Exp}(2\log\varphi\,H_\varphi)=F_{\mathrm{av}}^2\)}{Exp(2 log phi H_phi)=F_av^2}}

Let
\begin{equation}
 F_{\mathrm{av}}=\begin{pmatrix}0&1\\1&1\end{pmatrix},
 \qquad
 H_\varphi=\frac{2F_{\mathrm{av}}^2-3I}{2\varphi-1}.
 \label{eq:l9s102:euler-hphi}
\end{equation}
Then \(H_\varphi^2=I\) and
\begin{equation}
 \boxed{
 \operatorname{Exp}(2\log\varphi\,H_\varphi)=F_{\mathrm{av}}^2.}
 \label{eq:l9s102:euler-fav}
\end{equation}
In this coordination, \(\pi\) measures angular closure, \(\varphi\) the
expansive eigenvalue of memory, and \(e\) the operation that transforms an
additive generator into multiplicative transport.

\section{Ambivalence Principle: involution between \texorpdfstring{\(\pi/\varphi^2\)}{pi/phi^2} and \texorpdfstring{\(\varphi/\pi^2\)}{phi/pi^2} and nonadic accumulation of orientation}
\label{sec:l9s102:euler-ambivalencia}

\subsection{Action \texorpdfstring{\(\mathcal Q(x,y)=(x/y^2,y/x^2)\)}{Q(x,y)=(x/y^2,y/x^2)} on the areal and volumetric readers}

We define the quadratic reader
\begin{equation}
 \mathcal Q(x,y)
 =\left(\frac{x}{y^2},\frac{y}{x^2}\right).
 \label{eq:l9s102:euler-q}
\end{equation}
In the coordinates
\begin{equation}
 P=xy,
 \qquad O=\frac{x}{y},
 \label{eq:l9s102:euler-po}
\end{equation}
its action is
\begin{equation}
 P\circ\mathcal Q=P^{-1},
 \qquad
 O\circ\mathcal Q=O^3.
 \label{eq:l9s102:euler-q-action}
\end{equation}
The second iteration satisfies
\begin{equation}
 \boxed{
 P\circ\mathcal Q^2=P,
 \qquad
 O\circ\mathcal Q^2=O^9.}
 \label{eq:l9s102:euler-q2}
\end{equation}
The product returns and the orientation preserves a nonadic power of memory.

\subsection{Identity \texorpdfstring{\(\varphi^3(\pi/\varphi^2)^2(\varphi/\pi^2)=1\)}{phi^3(pi/phi^2)^2(phi/pi^2)=1} and typed bidirectionality}

For \((x,y)=(\pi,\varphi)\), the two orientations are
\begin{equation}
 \rho_A=\frac{\pi}{\varphi^2},
 \qquad
 \rho_E=\frac{\varphi}{\pi^2},
 \label{eq:l9s102:euler-rho-a-e}
\end{equation}
and obey
\begin{equation}
 \boxed{\varphi^3\rho_A^2\rho_E=1.}
 \label{eq:l9s102:euler-ambivalencia-identidad}
\end{equation}
The involution \((x,y)\mapsto(y,x)\) exchanges the areal and
electronic orientations of the same reader \(R(x,y)=x/y^2\).

\section{Exponential transformation of the parangular generator \texorpdfstring{\(A_{\mathrm{rad}}I+C^*_{\mathrm{rad}}H_\varphi\)}{A_rad I+C*_rad H_phi}}
\label{sec:l9s102:euler-accion-par-angular}

The action chart
\begin{equation}
 T_{-34}(u)=u\,10^{-34}\mathcal U_S
 \label{eq:l9s102:euler-t34}
\end{equation}
transports the angular identity
\(C^*_{\deg}=2(\eta_{\mathrm{ret}}+\alpha_{\mathrm{HMT}})\) to
\begin{equation}
 \boxed{
 T_{-34}(C^*_{\deg})
 =2\hbar_{\mathrm{ret}}
 +2T_{-34}(\alpha_{\angle,\deg}).}
 \label{eq:l9s102:euler-cstar-hbar}
\end{equation}
Each summand thus lies on the same dimensional line
\(\mathcal L_S\).

\subsection{Spectral projectors \texorpdfstring{\(P_\pm=(I\pm H_\varphi)/2\)}{P+/-=(I+/-H_phi)/2} and conjugate eigenvalues \texorpdfstring{\(q_\pm\)}{q+/-}}

Let
\begin{equation}
 B=A_{\mathrm{rad}}I+C^*_{\mathrm{rad}}H_\varphi,
 \qquad
 P_\pm=\frac{I\pm H_\varphi}{2}.
 \label{eq:l9s102:euler-b}
\end{equation}
Exponentiation transforms the additive generator into two conjugate multiplicative leaves:
\begin{equation}
 \boxed{
 e^{-B}=q_+P_++q_-P_-,
 \qquad
 q_\pm=e^{-A_{\mathrm{rad}}\mp C^*_{\mathrm{rad}}}.}
 \label{eq:l9s102:euler-exp-b}
\end{equation}

\subsection{Compatibility with the Parangular Publication and the Channels \texorpdfstring{\(90/120\)}{90/120}}

The cardinals \(90\) and \(120\) arise from the tritic incidence of the TPK.
The pair \(q_\pm\) does not generate them: it spectrally evaluates the operator
\(V_{90,120}\) on those previously constructed channels.

\section{Elliptic action boundary and hyperbolic interior of Hilbert--Beltrami--Klein}
\label{sec:l9s102:euler-elipse-klein}

\subsection{Conjugate coordinates of the action ellipse and law \texorpdfstring{\(\operatorname{Area}(\theta)=\hbar_{\mathrm{ret}}\theta\)}{Area(theta)=hbar_ret theta}}

The angular speed
\begin{equation}
 \eta=\operatorname{artanh}\!\left(\frac{C^*}{A}\right)
 \label{eq:l9s102:euler-rapidez}
\end{equation}
define the symplectic semiaxes
\begin{equation}
 a_S=\sqrt{2\hbar_{\mathrm{ret}}e^\eta},
 \qquad
 b_S=\sqrt{2\hbar_{\mathrm{ret}}e^{-\eta}}.
 \label{eq:l9s102:euler-semiejes}
\end{equation}
With
\begin{equation}
 Q=a_S\cos\theta,
 \qquad
 P=b_S\sin\theta,
 \label{eq:l9s102:euler-qp}
\end{equation}
the action enclosed by the arc satisfies
\begin{equation}
 \boxed{\operatorname{Area}(\theta)=\hbar_{\mathrm{ret}}\theta.}
 \label{eq:l9s102:euler-area-accion}
\end{equation}

\subsection{correspondence \texorpdfstring{\(\rho=\tan(\theta/2)=\tanh(u/2)\)}{rho=tan(theta/2)=tanh(u/2)} between borde periodico e interior hyperbolic}

The Gudermannian coordinate
\begin{equation}
 \rho=\tan\frac\theta2=\tanh\frac u2
 \label{eq:l9s102:euler-gudermann}
\end{equation}
coordinates the circular and hyperbolic charts. The boundary carries the
symplectic orbit; the convex interior domain carries the
Hilbert--Beltrami--Klein hyperbolic metric. The Gudermann coordinate realizes
their radial correspondence and preserves the distinction between metrics.

\section{Compatibility of the Euler operator closure with the Archimedean and exceptional publications}
\label{sec:l9s102:euler-doble-proyeccion}

\subsection{Double projection of the same enriched state and conservation of incidence}

Let \(X^{\mathrm{enr}}\) be the enriched state constructed by
APP--TRIT--TPK. Its two principal publications are
\begin{equation}
 X^{\mathrm{enr}}
 \xrightarrow{\ \mathcal P_{\mathrm{arq}}\ }\mathbb R,
 \qquad
 X^{\mathrm{enr}}
 \xrightarrow{\ \mathcal P_{\mathrm{exc}}\ }
 \mathrm{Witt}\to\mathrm{Golay}\to\mathrm{Leech}
 \to V^\natural\to\mathbb M\to J.
 \label{eq:l9s102:euler-doble-proyeccion}
\end{equation}
The first contracts compatible cylinders toward the Archimedean value; the
second preserves boundary incidence as far as the exceptional corridor.
The integer \(54\) appears in APP and in the modular coefficient of
\cref{eq:l9s102:euler-54}; the corridor \(4\to2\to6\to54\) subsequently transports
that incidence toward the Leech lattice, the vertex operator algebra,
and Moonshine. The action of the Monster resides in
\(V^\natural\); the digital and route genealogy remains in the domain
of TPK and reaches the corridor through the declared incidence maps.

\subsection{Sistema macrovectorial \texorpdfstring{\(\mathfrak E_{\mathrm{HMT}}=0\)}{E_HMT=0}: confluencia of the identities eulerianas in a unique operator}
\label{sec:l9s102:euler-macrovector}

The above identities are gathered in the typed application.
\begin{equation}
 \boxed{
 \mathfrak E_{\mathrm{HMT}}=
 \begin{pmatrix}
 (S_\times-S_+)-[q]t_3\\[1mm]
 \operatorname{Exp}(I)-eI\\
 \operatorname{Exp}(\pi J_{+1})+I\\
 \operatorname{Exp}(J_0)-I-J_0\\
 \operatorname{Exp}(2\log\varphi\,H_\varphi)-F_{\mathrm{av}}^2\\
 \varphi^3\rho_A^2\rho_E-1\\
 T_{-34}(C^*)-2\hbar_{\mathrm{ret}}
 -2T_{-34}(\alpha_{\angle,\deg})
 \end{pmatrix}=0.}
 \label{eq:l9s102:euler-macrovector}
\end{equation}
Each component retains its domain and its own proof. The aggregate
map exhibits the coordination of sum, product, elliptic closure, first
jet, hyperbolic expansion, ambivalence, and action.

The internal APP--TRIT--TPK composition already transports sum, product,
orientation, residue, quotient, carry, route, and memory to the
preceding publications. No absent HMT intertwiner is thereby posited. A
subsequent promotion can be formulated only through the specific typed square
required by a later representation---for example, complete group
equivariance as far as Witt---without downgrading the operator confluence
proved here.
